\chapter*{Abstract}

Automotive crash structures rely on aluminium components that see stress
states ranging from uniaxial tension to combined compression-shear during a
single impact event, yet most fracture-model calibration campaigns are built
almost entirely from tension-dominated specimens. This thesis addresses that
gap for cold-rolled AA6082 sheet, combining a rate-dependent necking study
with a damage-coupled fracture model calibrated jointly across tension- and
compression-dominated stress states.

Paper 1 characterises the rate dependence of diffuse necking onset in
AA6082-T6 sheet across four decades of strain rate, and shows that a modified
Considère criterion with a single rate-sensitivity parameter predicts the
observed onset strains to within 6\%. Paper 2 extends a local damage-coupled
constitutive model into the low- and negative-triaxiality regime using a
purpose-built compression-shear specimen, and shows that models calibrated
on tension data alone systematically over-predict ductility in this regime.

Together, the two papers give a calibration and validation route for
predicting ductile fracture in aluminium sheet across the full range of
stress states relevant to a crash event, rather than the tension-dominated
subset typically covered.
